\subsection{Internal Voltage and Arm Quantities}
\label{sec:converter-arm}

\subsubsection{Internal voltage and modulation index}

The upper and lower arm reactors of a phase act in parallel for the AC current, so the equivalent reactance per
phase is $X = \omega \cdot L_{arm} / 2$. With the valve-side phase voltage $U_{ph} = U_{valve} / \sqrt{3}$ as phase
reference, the AC current towards the grid and the internal voltage behind the arm reactors are
\begin{equation}
  \underline{I} = \frac{\left( P_{ac} + \mathrm{j} \cdot Q_{ac} \right)^{*}}{3 \cdot U_{ph}} , \qquad
  \underline{E} = U_{ph} + \mathrm{j} \cdot X \cdot \underline{I} , \qquad
  U_{conv} = \sqrt{3} \cdot |\underline{E}| .
\end{equation}
The arm reactance carries no active power; it absorbs $Q_{arm} = 3 \cdot X \cdot I^2$. The modulation index is the
peak phase voltage in per unit of half the DC voltage:
\begin{equation}
  m = \frac{2 \cdot \sqrt{2} \cdot U_{conv}}{\sqrt{3} \cdot U_{dc}} .
\end{equation}

\subsubsection{Arm voltages and currents}

Over one period, with $e(t) = \sqrt{2} \cdot |\underline{E}| \cdot \cos(\omega t)$ and the AC current $i_{ac}(t)$,
the arms of phase a carry
\begin{align}
  u_{upper}(t) &= \frac{U_{dc}}{2} - \left( e(t) + e_3(t) \right) , &
  i_{upper}(t) &= -\frac{I_{dc}}{3} - \frac{i_{ac}(t)}{2} , \\
  u_{lower}(t) &= \frac{U_{dc}}{2} + \left( e(t) + e_3(t) \right) , &
  i_{lower}(t) &= -\frac{I_{dc}}{3} + \frac{i_{ac}(t)}{2} ,
\end{align}
with the arm currents in their physical direction (upper arm towards the positive terminal) and $I_{dc} > 0$ into
the positive terminal (inverter). Each arm carries a third of the DC current and half of the AC current:
\begin{equation}
  I_{arm\_rms} = \sqrt{\left( \frac{I_{dc}}{3} \right)^2 + \left( \frac{I_{ac}}{2} \right)^2} , \qquad
  I_{arm\_peak} = \frac{|I_{dc}|}{3} + \frac{\sqrt{2} \cdot I_{ac}}{2} .
\end{equation}
The arm reactor losses are $P_{loss\_arm} = 6 \cdot R_{arm} \cdot I_{arm\_rms}^2$. The inserted submodules of an arm
follow its voltage, $n_{ins}(t) = u_{arm}(t) / u_{sm}$, bounded by 0 and $N$.

\subsubsection{Third harmonic injection}

With third harmonic injection the arm voltages contain the zero-sequence component
$e_3(t) = -\tfrac{\sqrt{2} \cdot |\underline{E}|}{6} \cdot \cos(3 \omega t)$. It cancels between the phases and does
not change the AC side; it lowers the peak of $e + e_3$ to $\sqrt{3}/2$ of the fundamental peak. The peak factor is
$k_3 = 1$ without and $k_3 = \sqrt{3}/2$ with injection.

\subsubsection{Limits}

A half-bridge arm cannot produce a negative voltage, and its voltage cannot exceed the sum of its capacitor
voltages:
\begin{equation}
  k_3 \cdot m \le 1 , \qquad
  U_{arm\_max} = \frac{U_{dc}}{2} \cdot \left( 1 + k_3 \cdot m \right) \le N \cdot u_{sm} .
\end{equation}
The highest modulation index at a given DC voltage is therefore
$m_{max} = \min\left( 1, \, 2 \cdot N \cdot u_{sm} / U_{dc} - 1 \right) / k_3$. An exceeded limit is a warning of
the result; the calculation is not stopped. The reactive power the converter can provide at a given active power is
the circle $|\underline{E}| \le E_{max}$ in the $PQ$ plane, with centre $Q = -U_{valve}^2 / X$ and radius
$U_{valve} \cdot E_{max} / X$, $E_{max} = m_{max} \cdot \sqrt{3} \cdot U_{dc} / (2 \cdot \sqrt{2})$ line-to-line.

\subsubsection{DC side}

The DC power is the AC power plus the losses, and the DC current follows from the DC voltage:
\begin{equation}
  P_{dc} = P_{ac} + P_{loss\_cv} , \qquad I_{dc} = \frac{P_{dc}}{U_{dc}} ,
\end{equation}
both in the convention of Section~\ref{sec:conventions} ($P_{dc} < 0$ for a rectifier). The converter can also be
calculated with $P_{dc}$ or $I_{dc}$ given; then $P_{ac} = P_{dc} - P_{loss\_cv}(P_{ac})$ is iterated to a change
below $10^{-10}$ of $|P_{dc}|$.
